\chapter{L'architecture Transformer}
\label{ch:transformers}

\begin{quote}
\emph{<< Attention is all you need. >> --- Vaswani et al., 2017}
\end{quote}

% ============================================================
\section{Des RNN aux Transformers}

Les r\'eseaux r\'ecurrents (RNN, LSTM) traitent les s\'equences token par token, de gauche \`a droite. Cela cr\'ee deux probl\`emes :
\begin{enumerate}
  \item \textbf{Goulet d'\'etranglement s\'equentiel.} Le token $t$ doit attendre que les tokens $1, \ldots, t-1$ soient trait\'es. Pas de parall\'elisme.
  \item \textbf{Contexte qui s'estompe.} L'information des premiers tokens s'efface \`a mesure que la s\'equence grandit.
\end{enumerate}

Le Transformer r\'esout les deux probl\`emes avec l'\emph{attention} : chaque token peut regarder chaque autre token directement, en parall\`ele.

% ============================================================
\section{Attention par produit scalaire \`a l'\'echelle}

\'Etant donn\'ee une s\'equence de $n$ tokens, chacun repr\'esent\'e par un vecteur, on calcule trois matrices :
\begin{itemize}
  \item $\mathbf{Q}$ (queries), $\mathbf{K}$ (keys), $\mathbf{V}$ (values), chacune de forme $(n \times d_k)$.
\end{itemize}

\[
\text{Attention}(\mathbf{Q}, \mathbf{K}, \mathbf{V}) = \text{softmax}\!\left(\frac{\mathbf{Q}\mathbf{K}^\top}{\sqrt{d_k}}\right)\mathbf{V}
\]

Le facteur $\sqrt{d_k}$ emp\^eche les produits scalaires de devenir trop grands, ce qui pousserait le softmax dans des r\'egions \`a gradients qui s'\'evanouissent.

\begin{minted}{python}
import torch
import torch.nn.functional as F

def scaled_dot_product_attention(Q, K, V, mask=None):
    """Compute scaled dot-product attention."""
    d_k = Q.size(-1)
    scores = torch.matmul(Q, K.transpose(-2, -1)) / (d_k ** 0.5)
    if mask is not None:
        scores = scores.masked_fill(mask == 0, float("-inf"))
    weights = F.softmax(scores, dim=-1)
    return torch.matmul(weights, V), weights

# Exemple : 1 batch, 4 tokens, dimension de plongement 8
Q = torch.randn(1, 4, 8)
K = torch.randn(1, 4, 8)
V = torch.randn(1, 4, 8)
output, attn_weights = scaled_dot_product_attention(Q, K, V)
print(f"Output shape: {output.shape}")       # (1, 4, 8)
print(f"Attention weights shape: {attn_weights.shape}")  # (1, 4, 4)
\end{minted}

\begin{aitip}
La matrice de poids d'attention est $(n \times n)$. L'entr\'ee $(i, j)$ vous dit \`a quel point le token $i$ << regarde >> le token $j$. Visualiser cette matrice r\'ev\`ele ce sur quoi le mod\`ele se concentre.
\end{aitip}

% ============================================================
\section{Attention multi-t\^etes}

Au lieu d'une seule fonction d'attention, les Transformers utilisent $h$ << t\^etes >> parall\`eles, chacune avec ses propres projections apprises $W_Q^{(i)}, W_K^{(i)}, W_V^{(i)}$ :

\[
\text{MultiHead}(\mathbf{Q}, \mathbf{K}, \mathbf{V}) = \text{Concat}(\text{head}_1, \ldots, \text{head}_h)\,\mathbf{W}^O
\]

Chaque t\^ete peut apprendre un type de relation diff\'erent : syntaxique, s\'emantique, positionnel, etc.

\begin{minted}{python}
import torch.nn as nn

class MultiHeadAttention(nn.Module):
    def __init__(self, d_model=512, num_heads=8):
        super().__init__()
        assert d_model % num_heads == 0
        self.d_k = d_model // num_heads
        self.num_heads = num_heads
        self.W_q = nn.Linear(d_model, d_model)
        self.W_k = nn.Linear(d_model, d_model)
        self.W_v = nn.Linear(d_model, d_model)
        self.W_o = nn.Linear(d_model, d_model)

    def forward(self, Q, K, V, mask=None):
        batch = Q.size(0)
        # Projection puis reshape : (batch, seq, d_model) -> (batch, heads, seq, d_k)
        Q = self.W_q(Q).view(batch, -1, self.num_heads, self.d_k).transpose(1, 2)
        K = self.W_k(K).view(batch, -1, self.num_heads, self.d_k).transpose(1, 2)
        V = self.W_v(V).view(batch, -1, self.num_heads, self.d_k).transpose(1, 2)
        # Attention par t\^ete
        out, weights = scaled_dot_product_attention(Q, K, V, mask)
        # Concat\'enation des t\^etes
        out = out.transpose(1, 2).contiguous().view(batch, -1, self.num_heads * self.d_k)
        return self.W_o(out)

mha = MultiHeadAttention(d_model=64, num_heads=4)
x = torch.randn(2, 10, 64)  # batch=2, seq=10, dim=64
print(f"Output: {mha(x, x, x).shape}")  # (2, 10, 64)
\end{minted}

% ============================================================
\section{Encodage positionnel}

L'attention est invariante par permutation --- elle ne conna\^it pas l'ordre des tokens. Les encodages positionnels injectent l'information de position :

\[
PE_{(pos, 2i)} = \sin\!\left(\frac{pos}{10000^{2i/d}}\right), \quad
PE_{(pos, 2i+1)} = \cos\!\left(\frac{pos}{10000^{2i/d}}\right)
\]

\begin{minted}{python}
import numpy as np
import matplotlib.pyplot as plt

def positional_encoding(max_len, d_model):
    pe = np.zeros((max_len, d_model))
    position = np.arange(max_len)[:, np.newaxis]
    div_term = np.exp(np.arange(0, d_model, 2) * -(np.log(10000.0) / d_model))
    pe[:, 0::2] = np.sin(position * div_term)
    pe[:, 1::2] = np.cos(position * div_term)
    return pe

pe = positional_encoding(100, 64)
plt.figure(figsize=(10, 4))
plt.imshow(pe, aspect="auto", cmap="RdBu")
plt.xlabel("Embedding dimension")
plt.ylabel("Position")
plt.title("Sinusoidal positional encoding")
plt.colorbar()
plt.tight_layout()
plt.savefig("positional_encoding.png", dpi=150)
plt.show()
\end{minted}

\begin{aitip}
Les mod\`eles modernes (GPT, Llama) utilisent les \emph{Rotary Position Embeddings} (RoPE) plut\^ot que les encodages sinuso\"idaux. RoPE code la position relative en faisant tourner les vecteurs query et key, ce qui passe mieux \`a l'\'echelle des contextes longs.
\end{aitip}

% ============================================================
\section{Architecture encodeur-d\'ecodeur}

Le Transformer original a deux piles :
\begin{itemize}
  \item \textbf{Encodeur :} traite la s\'equence d'entr\'ee avec auto-attention bidirectionnelle. Chaque token peut voir tous les autres tokens.
  \item \textbf{D\'ecodeur :} g\'en\`ere la s\'equence de sortie avec auto-attention causale (masqu\'ee). Chaque token ne peut voir que les tokens pr\'ec\'edents. Il attribue \'egalement de l'attention \`a la sortie de l'encodeur via la cross-attention.
\end{itemize}

\begin{center}
\begin{tikzpicture}[
  block/.style={rectangle, draw, fill=aiblue!15, minimum width=3cm, minimum height=0.8cm, rounded corners},
  arrow/.style={-{Stealth}, thick}
]
  \node[block] (emb_enc) at (0, 0) {Plongement entr\'ee + PE};
  \node[block] (sa_enc) at (0, 1.2) {Auto-attention multi-t\^etes};
  \node[block] (ff_enc) at (0, 2.4) {R\'eseau Feed-Forward};
  \node[above] at (0, 3.1) {\textbf{Encodeur} ($\times N$)};

  \node[block] (emb_dec) at (5.5, 0) {Plongement sortie + PE};
  \node[block] (sa_dec) at (5.5, 1.2) {Auto-attention masqu\'ee};
  \node[block] (ca_dec) at (5.5, 2.4) {Cross-attention};
  \node[block] (ff_dec) at (5.5, 3.6) {R\'eseau Feed-Forward};
  \node[block] (linear) at (5.5, 4.8) {Lin\'eaire + Softmax};
  \node[above] at (5.5, 5.5) {\textbf{D\'ecodeur} ($\times N$)};

  \draw[arrow] (emb_enc) -- (sa_enc);
  \draw[arrow] (sa_enc) -- (ff_enc);
  \draw[arrow] (emb_dec) -- (sa_dec);
  \draw[arrow] (sa_dec) -- (ca_dec);
  \draw[arrow] (ca_dec) -- (ff_dec);
  \draw[arrow] (ff_dec) -- (linear);
  \draw[arrow, dashed] (ff_enc.east) -- ++(1,0) |- (ca_dec.west);
\end{tikzpicture}
\end{center}

% ============================================================
\section{Variantes en pratique}

\begin{longtable}{p{3.5cm}p{3.5cm}p{5cm}}
\toprule
\textbf{Architecture} & \textbf{Mod\`eles} & \textbf{Cas d'usage} \\
\midrule
Encodeur seul & BERT, DistilBERT, RoBERTa & Classification, NER, plongements \\
D\'ecodeur seul & GPT-2, GPT-4, Llama, Phi & G\'en\'eration de texte, chatbots, code \\
Encodeur-d\'ecodeur & T5, BART, Flan-T5 & Traduction, r\'esum\'e \\
\bottomrule
\end{longtable}

% ============================================================
\section{Visualiser l'attention avec DistilBERT}

\begin{minted}{python}
from transformers import AutoTokenizer, AutoModel
import torch
import matplotlib.pyplot as plt
import seaborn as sns

model_name = "distilbert-base-uncased"
tokenizer = AutoTokenizer.from_pretrained(model_name)
model = AutoModel.from_pretrained(model_name, output_attentions=True)

text = "The cat sat on the mat because it was tired"
inputs = tokenizer(text, return_tensors="pt")
with torch.no_grad():
    outputs = model(**inputs)

# outputs.attentions est un tuple : un tenseur par couche
# Forme de chaque tenseur : (batch, heads, seq_len, seq_len)
attn = outputs.attentions[-1][0]  # derni\`ere couche, premier batch
tokens_list = tokenizer.convert_ids_to_tokens(inputs["input_ids"][0])

fig, axes = plt.subplots(1, 4, figsize=(20, 5))
for i, ax in enumerate(axes):
    sns.heatmap(attn[i].numpy(), xticklabels=tokens_list,
                yticklabels=tokens_list, ax=ax, cmap="Blues")
    ax.set_title(f"Head {i}")
plt.tight_layout()
plt.savefig("attention_heads.png", dpi=150)
plt.show()
\end{minted}

\begin{exercise}
\begin{enumerate}
  \item Impl\'ementez la fonction \texttt{scaled\_dot\_product\_attention} \`a partir de z\'ero (sans l'API d'attention PyTorch). V\'erifiez qu'elle reproduit la sortie de \texttt{torch.nn.functional.scaled\_dot\_product\_attention}.
  \item Visualisez les t\^etes d'attention pour la phrase << The bank by the river was steep. >>. Quelles t\^etes attribuent leur attention \`a quels mots ? Y a-t-il des t\^etes qui capturent le lien entre << bank >> et << river >> ?
  \item Modifiez la fonction d'encodage positionnel pour utiliser des plongements appris plut\^ot que sinuso\"idaux. Entra\^inez un petit mod\`ele et comparez.
  \item Calculez le nombre de param\`etres d'une couche d'encodeur Transformer avec $d_{\text{model}}=512$ et $h=8$. D\'etaillez votre calcul.
\end{enumerate}
\end{exercise}

\begin{ethicsbox}
Les poids d'attention sont parfois interpr\'et\'es comme des << explications >> du comportement du mod\`ele. C'est trompeur. L'attention montre o\`u le mod\`ele \emph{regarde}, pas n\'ecessairement ce qu'il \emph{utilise} pour la pr\'ediction finale. La visualisation d'attention est un outil de diagnostic, pas une explication fid\`ele.
\end{ethicsbox}

% ============================================================
\section{R\'esum\'e du chapitre}

\begin{itemize}
  \item Le Transformer remplace la r\'ecurrence par l'auto-attention, ce qui permet le parall\'elisme complet.
  \item L'attention par produit scalaire \`a l'\'echelle calcule une somme pond\'er\'ee des valeurs, avec des poids issus de la similarit\'e query-key.
  \item L'attention multi-t\^etes permet au mod\`ele d'attribuer son attention \`a diff\'erents types de relations simultan\'ement.
  \item Les encodages positionnels (sinuso\"idaux ou appris) injectent l'ordre s\'equentiel dans le m\'ecanisme d'attention.
  \item Encodeur seul (BERT), d\'ecodeur seul (GPT), et encodeur-d\'ecodeur (T5) sont les trois principales variantes de Transformer.
  \item La visualisation d'attention est utile pour le d\'ebogage, mais ne doit pas \^etre sur-interpr\'et\'ee comme une explication du mod\`ele.
\end{itemize}
